% Input immediately after mean_field.tex, before spectral_obstruction.tex.
% A scalar refinement; the original mf:257 family remains unchanged.

\subsection{A refinement to limiting condition number 227}
\label{rf:subsection}

A different choice of the two scalar parameters improves the conditioning
bound using the same spherical law, ordering, and sector estimates.

\begin{theorem}[An eventual family with condition number below \(228\)]
\label{rf:228}
There is an explicit algebraic sequence of real positive-definite
correlation matrices \(\widetilde E_N\), defined in every sufficiently
large order, such that
\[
 \lim_{N\to\infty}
 \frac{D_1(\widetilde E_N)}{N(N-1)\per\widetilde E_N}
 <-\frac1{40000},\qquad
 |(\widetilde E_N)_{ij}|\le\frac{904}{N}\quad(i\ne j).
\]
Its least eigenvalue tends to one, and its condition number tends to
\(227\), so it is eventually below \(228\). Before correlation
normalization it is a mean-field perturbation with latent rank four.
No explicit sufficient order is asserted.
\end{theorem}
\begin{proof}
Use the construction of Theorem~\ref{mf:257}, replacing its parameters
by \(\tau=9999/10000\) and \(t=904\). Write
\(m(k)=1-\log(1+k)/k\) and \(t(k)=2k/m(k)\).
The elementary bounds \(5/2<e<11/4\),
\((11/4)^6<445\), and \((5/2)^7>447\) give
\(6<\log445<\log447<7\). Consequently
\[
 t(444)<\frac{2\cdot444^2}{437}
       =\frac{394272}{437}<904,
 \qquad
 t(446)>\frac{2\cdot446^2}{440}
       =\frac{49729}{55}>904.
\]
Strict monotonicity of \(t(k)\) proves \(444<k<446\).
With \(K=1+k\) and \(\delta=1-m\), this implies
\[
 \frac6{446}<\delta<\frac7{444},\qquad
 \frac{2109}{100}<\sqrt K<\frac{423}{20}.
\]
The square-root inequalities follow by squaring:
\((2109/100)^2<445\) and \((423/20)^2>447\).
Using \(\pi>31415/10000\) in Lemma~\ref{mf:four-bounds} gives
\begin{align*}
 \Gamma_k
 &>\frac{31415/40000-2(423/20)/444}{1-6/446}
   =\frac{45552433}{65120000}=:g>0,\\
 \rho_k
 &>\frac12-\frac7{888}-\frac{100}{2209}
   =\frac{876533}{1961592}=:r>0.
\end{align*}
As in \eqref{mf:four-rational-margin}, use \(1/m>1\) and
\(\delta^2\ge\tau^2\delta^2\) in the full score. Exact rational
arithmetic yields
\begin{align}
 \mathcal R_*
 &>\left(\frac{9999}{10000}\right)^2
  \left[2g^2+8\frac3{223}r^2+
                     \left(\frac3{223}\right)^2\right]\notag\\
 &=\frac{425269160832679340834525001397641}
         {425220873828711680000000000000000}
 >1+\frac1{10000}.\label{rf:rational-margin}
\end{align}
Theorem~\ref{mf:criterion} and \eqref{mf:ratio} give the stated
derivative limit.

Use the same algebraic sphere grid, padding, and sorting as in
Theorem~\ref{mf:257}, with the new \(\tau\), and set
\(\widetilde A_N=I_N+(904/N)V_NV_N^T\). The row norms are at most
one, and the limiting covariance is
\(\operatorname{diag}(1/4,1/4,\tau^2/4,\tau^2/4)\).
The condition number therefore tends to \(1+904/4=227\).
Correlation normalization preserves the normalized derivative, changes
the operator by \(o(1)\), and gives the off-diagonal bound \(904/N\)
and the least-eigenvalue limit as before.
\end{proof}

The family in Theorem~\ref{mf:257} retains the larger proved derivative
margin \(1/1000\). The present theorem improves its condition-number
bound with a smaller proved margin; it is a scalar refinement of the
same analytic construction.
